\section{Selección de Datos}
Para el desarrollo del presente proyecto se seleccionaron diferentes conjuntos de datos provenientes de fuentes oficiales, con el fin de construir un análisis integral que permita relacionar variables educativas, socioeconómicas y de conectividad a nivel territorial.
La selección de estos datasets responde directamente al objetivo del negocio, el cual busca analizar la relación entre el acceso a internet, las condiciones socioeconómicas y el desempeño académico en Colombia.
\subsection{Resultados Saber 11 (ICFES)}
Este conjunto de datos contiene los resultados individuales de las pruebas Saber 11, incluyendo puntajes por áreas y variables socioeconómicas de los estudiantes.
Su principal relevancia radica en que permite analizar directamente el desempeño académico, siendo esta la variable dependiente del proyecto. Adicionalmente, incluye información sobre acceso a internet y disponibilidad de computador en el hogar, lo que permite establecer relaciones directas entre condiciones tecnológicas y resultados educativos.
\subsection{Acceso a Internet Fijo por Municipio (MinTIC)}
Este dataset proporciona información sobre la cantidad de accesos a internet fijo, desagregados por tecnología y segmento a nivel municipal.
Su uso permite caracterizar el nivel de conectividad en los territorios, lo que resulta fundamental para analizar la brecha digital y su posible impacto en el desempeño educativo. Asimismo, permite comparar municipios con distintos niveles de penetración de internet y evaluar diferencias en sus resultados académicos.
\subsection{Datos de Sisbén / Condiciones Socioeconómicas (DNP)}
Los datos del Sisbén permiten identificar condiciones de vulnerabilidad de la población, incluyendo variables relacionadas con nivel socioeconómico, acceso a servicios y condiciones de vida.
Estos datos son fundamentales para complementar el análisis, ya que permiten evaluar el impacto de la pobreza y otras privaciones sobre el desempeño académico, así como identificar si la conectividad por sí sola es suficiente para mejorar los resultados educativos.
\subsection{Datos de Educación por Municipio (MEN)}
Este conjunto de datos incluye información relacionada con matrícula, cobertura y características del sistema educativo a nivel municipal.
Su inclusión permite contextualizar los resultados académicos dentro de la estructura del sistema educativo, facilitando el análisis de variables como cobertura escolar y distribución territorial de la educación.
El dataset contiene 15,707 registros sin duplicados, distribuidos en 1,124 municipios de 36 departamentos, con cobertura temporal de 2011 a 2024 (14 vigencias). Cada registro corresponde a un municipio-año y agrupa 41 variables, de las cuales 34 son indicadores numéricos que cubren cinco dimensiones educativas: cobertura (neta y bruta por nivel), tasas de deserción intra-anual, aprobación, reprobación y repitencia — cada una desagregada por nivel educativo (transición, primaria, secundaria y media) — y el porcentaje de sedes oficiales conectadas a internet. Los identificadores territoriales incluyen código y nombre del municipio, departamento y Entidad Territorial Certificada (ETC).